% ============================================================
% SCIENTIFIC PAPER (ban co dong) - Laptop Selection using AHP-TOPSIS
% Bien dich: xelatex paper.tex -> biber paper -> xelatex paper x2
% Dung chung references.bib + figures/ voi Report (main.tex).
% Single-column journal style (article), APA narrative citation.
% ============================================================
\documentclass[11pt,a4paper]{article}

\usepackage{fontspec}
\usepackage{polyglossia}
\setmainlanguage{english}
\setmainfont{Times New Roman}

\usepackage[margin=2.5cm]{geometry}
\usepackage{booktabs}
\usepackage{graphicx}
\graphicspath{{figures/}}
\usepackage{amsmath,amssymb,mathtools}
\usepackage{microtype}
\usepackage{setspace}
\onehalfspacing

\usepackage[style=apa,backend=biber]{biblatex}
\addbibresource{references.bib}

\usepackage{hyperref}
\usepackage{cleveref}

\title{A Hybrid AHP--TOPSIS Approach to Laptop Selection\\for Office Workers in Vietnam}
\author{Group [X] --- Logistics Management\\University of Finance --- Marketing}
\date{\today}

\begin{document}
\maketitle

\begin{abstract}
% TODO Phase 6: Abstract ~200-250 tu, khuon gap -> method -> result -> implication (mau Maghsoudi 2026).
\noindent\textit{[Abstract to be written in Phase 6 --- gap $\rightarrow$ method (AHP--TOPSIS, three stages) $\rightarrow$ results (weights, ranking) $\rightarrow$ implication.]}

\vspace{0.5em}
\noindent\textbf{Keywords:} Multi-Criteria Decision-Making; AHP; TOPSIS; laptop selection; office workers; Vietnam.
\end{abstract}

% ============================================================
\section{Introduction}\label{sec:p-intro}
% Ban co dong cua Chapter 1 (Report). Giu citation that; nen ~2-3 doan.

The laptop has become a basic working tool across most occupations, and its consumption has
expanded steadily. \textcite{sonmez_cakir_determination_2020} report that 161.6 million laptops
were sold in 2017, and \textcite{gaur_application_2023} record a rise to 218 million units in
2020 and 225 million in 2021. As the catalogue fragments into many brands, models, and
configurations, buyers must compare technical attributes expressed on different scales that
frequently conflict, which places the choice within the scope of Multi-Criteria Decision-Making
(MCDM). Among MCDM techniques, the Analytic Hierarchy Process (AHP) and the Technique for Order
of Preference by Similarity to Ideal Solution (TOPSIS) are among the most widely used
\parencite{basilio_systematic_2022,zyoud_bibliometric-based_2017}, and their combination recurs
in applied research \parencite{tighnavard_balasbaneh_systematic_2025}.

% TODO Phase 6: doan 2 — problem statement + gap + twofold contribution (nen tu Ch.1 muc 1.2-1.3).
% TODO Phase 6: doan cuoi — 3 RQ + cau "The remainder of this paper is organised as follows".
\textit{[Phase 6 — condense problem statement, gap, contributions, and research questions; see \texttt{SUON-NOI-DUNG.md} \S P1.]}

% ============================================================
\section{Literature Review}\label{sec:p-lit}
% TODO Phase 6: nen ~1-1.5 trang — tong quan MCDM (phan loai 5 nhom Villalba) + dinh vi
% AHP-TOPSIS + gap. Co the giu bang tong hop rut gon (Source | Product | Criteria | Method).
\textit{[Phase 6 — see \texttt{SUON-NOI-DUNG.md} \S P2.]}

% ============================================================
\section{Methodology}\label{sec:p-method}
% Method da chot (khong phu thuoc Phase 4-5) — co the giu day du. Nen 3 tieu muc.

\subsection{Research Framework and Criteria}\label{sec:p-framework}
% TODO Phase 6: so do 3 giai doan + bang 14 tieu chi/6 nhom (rut gon tu Report Table).
\textit{[Phase 6 — three-stage framework and the fourteen-criteria table; see \texttt{SUON-NOI-DUNG.md} \S P3.]}

\subsection{AHP for Criteria Weighting}\label{sec:p-ahp}
Criteria weights are derived through AHP. Expert pairwise judgements on the Saaty $1$--$9$ scale
form a reciprocal matrix $A=[a_{ij}]$; the weight of each criterion is the normalised row average
\begin{equation}\label{eq:p-weight}
  w_i = \frac{1}{n}\sum_{j=1}^{n}\frac{a_{ij}}{\sum_{k=1}^{n}a_{kj}},
\end{equation}
and consistency is verified through $CR = CI/RI$ with
\begin{equation}\label{eq:p-cr}
  CI = \frac{\lambda_{\max}-n}{\,n-1\,}, \qquad CR \le 0.10.
\end{equation}

\subsection{TOPSIS for Ranking}\label{sec:p-topsis}
Given the weights, the candidate laptops are ranked with TOPSIS. The decision matrix is
vector-normalised, $r_{ij}=x_{ij}/\sqrt{\sum_i x_{ij}^2}$, and weighted, $v_{ij}=w_j r_{ij}$.
After forming the positive and negative ideal solutions $A^{+}$ and $A^{-}$, the separation
measures $S_i^{+}$ and $S_i^{-}$ yield the closeness coefficient
\begin{equation}\label{eq:p-closeness}
  C_i = \frac{S_i^{-}}{\,S_i^{+}+S_i^{-}\,},
\end{equation}
and alternatives are ranked in descending order of $C_i$.

% ============================================================
\section{Results}\label{sec:p-results}
% TODO Phase 4-5: bang trong so + CR; ma tran quyet dinh; bang xep hang C_i; hinh do nhay.
\textit{[Phase 4-5 output — see \texttt{SUON-NOI-DUNG.md} \S P4.]}

% ============================================================
\section{Discussion}\label{sec:p-discussion}
% TODO Phase 6: phat hien ly thuyet quoc te vs thuc tien VN; ham y; doi chieu nghien cuu truoc.
\textit{[Phase 6 — see \texttt{SUON-NOI-DUNG.md} \S P5.]}

% ============================================================
\section{Conclusion}\label{sec:p-conclusion}
% TODO Phase 6: tra loi 3 RQ + dong gop + han che + huong phat trien (gop gon).
\textit{[Phase 6 — see \texttt{SUON-NOI-DUNG.md} \S P6.]}

\printbibliography

\end{document}
